\documentclass[a4paper,11pt]{article}
\usepackage{iftex}
\ifPDFTeX
  \usepackage[utf8]{inputenc}
  \usepackage[T1]{fontenc}
  \usepackage{lmodern}
\else
  \usepackage{fontspec}
  \setmainfont{lmroman10-regular.otf}[
    BoldFont=lmroman10-bold.otf,
    ItalicFont=lmroman10-italic.otf,
    BoldItalicFont=lmroman10-bolditalic.otf]
  \setmonofont{lmmono10-regular.otf}
\fi
\usepackage[croatian]{babel}
\usepackage[margin=24mm]{geometry}
\usepackage{enumitem}
\usepackage[hidelinks]{hyperref}
\setlength{\parindent}{0pt}
\setlength{\parskip}{6pt}
\setlist[itemize]{nosep,leftmargin=16pt,topsep=3pt}
\setlist[enumerate]{nosep,leftmargin=18pt,topsep=3pt}
\hypersetup{pdftitle={Data Monitoring Platform - Izvješće 2},pdfauthor={Leonard Lečić i Dario Golubović}}
\AtEndDocument{\clearpage\ifnum\value{page}=3\relax\else
\PackageError{KT2}{Izvjesce mora imati dvije stranice}{Prilagodite duljinu sadrzaja.}\fi}
\begin{document}
\begin{center}
{\Large\bfseries Data Monitoring Platform}\\[5pt]
{\large Izvješće 2 -- Druga kontrolna točka}\\[7pt]
Leonard Lečić i Dario Golubović\\
Skriptni jezici \quad | \quad 25. studenoga 2026.
\end{center}



\section*{1. Napredak u drugoj kontrolnoj točki}
Napravili smo Python program koji prikuplja stanje Bajs stanica u
Zagrebu i sprema podatke u lokalnu bazu. Dohvat radi automatski
svakih pet minuta, a može se pokrenuti i ručno. Tako smo počeli
graditi povijest iz koje ćemo kasnije vidjeti koje stanice najčešće
ostaju bez bicikala i u koje doba dana.

\section*{2. Automatsko prikupljanje}
Podatke preuzimamo iz javnog Nextbike API-ja. Za svaku stanicu
spremamo njezin broj i naziv, broj dostupnih bicikala, broj slobodnih
mjesta kad je dostupan, oznaku održavanja i vrijeme dohvata.
Prije spremanja provjeravamo jesu li potrebni podaci prisutni,
jesu li odgovarajućeg tipa i je li broj bicikala valjan.

Naredba \texttt{collect} pokreće jedan dohvat, a \texttt{watch}
ponavlja isti postupak u zadanom razmaku. Početni razmak je pet
minuta i može se promijeniti u konfiguracijskoj datoteci.
U dnevnik se zapisuju uspješni i neuspjeli pokušaji. Ako dohvat
ne uspije, program ponovno pokušava u sljedećem ciklusu.

Na računalu koje skuplja podatke postavili smo Windows zadatak
koji pokreće program pri prijavi korisnika. Program radi u pozadini
i traži da računalo tijekom prikupljanja ostane budno.
Zadatak dopušta jedno pokretanje programa u isto vrijeme.

\section*{3. Spremanje i pregled podataka}
Za pohranu koristimo \textbf{SQLite}. Svaki redak u bazi opisuje
jednu stanicu u trenutku dohvata. Vrijeme spremamo u UTC-u, a
identifikator stanice i vrijeme snimke zajedno sprječavaju
dvostruko spremanje istog zapisa. Podaci iz jednog dohvata
spremaju se zajedno.

Naredba \texttt{status} prikazuje vrijeme zadnje snimke, broj
aktivnih stanica, koliko ih je bez bicikala i koliko je stanica
na održavanju. Dodali smo i uvoz ranije spremljenog JSON odgovora
uz njegovo izvorno vrijeme dohvata. To koristimo za provjere i
ponavljanje istog postupka bez pristupa API-ju.

Program je napisan u \textbf{Pythonu 3.12+}, uz standardne module
za mrežne zahtjeve, JSON, konfiguraciju, bazu i dnevnike rada.
Windows zadatak postavljamo \textbf{PowerShell} skriptom.
Za testiranje koristimo \textbf{pytest}, a za provjeru koda
\textbf{Ruff}.

\newpage
\section*{4. Prikupljeni podaci}
Prikupljanje smo pokrenuli 12. studenoga 2026. u 20:01.
Do 25. studenoga u 18:17 prikupili smo \textbf{3.683 snimke}
s po 200 stanica, odnosno \textbf{736.600 zapisa}.
Baza je tada zauzimala približno 93~MB.

U zadnjoj snimci bilo je 155 aktivnih stanica, od kojih je deset
bilo bez dostupnih bicikala. Još 45 stanica bilo je označeno za
održavanje. Ovu oznaku uzet ćemo u obzir pri analizi dostupnosti.

Za izvješće smo sačuvali kopiju baze i provjerili njezinu
ispravnost. Prikupljanje se nastavlja kako bismo dobili više
podataka za daljnju analizu. Kod, testovi i izvješća nalaze se
u Git repozitoriju, a baza i dnevnici čuvaju se lokalno.

\section*{5. Provjere i poteškoće}
Prošlo je svih \textbf{26 testova}, a kod je prošao i Ruff provjeru.
Provjerili smo konfiguraciju, obradu odgovora, uvoz spremljenih
podataka i prikaz zadnje snimke. Provjerili smo i prepoznavanje
već spremljenih snimki te nastavak prikupljanja nakon mrežne
pogreške.

Prvi veći problem pojavio se 13. studenoga, kada je računalo
prešlo u mirovanje. Ponovno smo pokrenuli program i dodali
zaštitu od automatskog mirovanja tijekom prikupljanja.
Dogodili su se i pojedinačni problemi s mrežom ili čekanjem
odgovora API-ja. Program ih je zabilježio i nastavio raditi
u sljedećem ciklusu. Takvi prekidi vidljivi su kao praznine
između spremljenih snimki.

\section*{6. Plan za nastavak rada}
Za KT3, do 18. prosinca, planiramo napraviti web stranicu s
kartom stanica. Klikom na stanicu otvorit će se njezina povijest,
graf po satima i osnovna statistika. Dodat ćemo popis stanica
koje najčešće ostaju bez bicikala te mogućnost odabira favorita
i uobičajenog vremena korištenja.

Za KT4, do 22. siječnja, planiramo povezati podatke s vremenskim
API-jem Open-Meteo i usporediti dostupnost bicikala u kišnim i
suhim terminima. Na temelju rezultata i korisničkih postavki
dodat ćemo obavijesti e-poštom. Na kraju ćemo provjeriti cijeli
postupak, ispraviti preostale greške i dovršiti upute.
Posao ćemo podijeliti na analizu i prikaz te povezivanje s
vremenom i obavijestima, uz zajedničku provjeru aplikacije.

\end{document}
